% titel: Federpendel (Beispiel)
% skript: Muster ohne Skriptbezug
% stand: Beispiel, darf gelöscht werden, sobald echte Aufgaben da sind

% Teil I: nur Werkzeuge. Kein Lösungsweg, keine Zahlen dieser Aufgabe.
% Der Build setzt die Abschnittsüberschrift und den Lösungsverweis.

Ein Körper der Masse $m$ hängt reibungsfrei an einer Feder mit der
Federkonstante $k$ und wird ein Stück aus der Ruhelage gezogen. Gefragt ist
nach dem zeitlichen Verlauf der Bewegung und nach Periodendauer, Frequenz und
Grenzgeschwindigkeit.

\begin{werkzeugkasten}[Werkzeuge]
\textbf{Zweites Newtonsches Gesetz} in einer Dimension, hier genügt der Betrag
mit Vorzeichen in Auslenkungsrichtung:
\begin{equation}
  m\,\ddot{x} = \sum_i F_i .
\end{equation}

\textbf{Hookesche Kraft} der idealen Feder, immer zurück zur Ruhelage zeigend:
\begin{equation}
  F_{\text{Feder}} = -k\,x, \qquad [k] = \si{\newton\per\metre} = \si{\kilo\gram\per\second\squared}.
\end{equation}

\textbf{Bewegungsgleichung des harmonischen Oszillators} nach Einsetzen und
Umstellen:
\begin{equation}
  \ddot{x} + \omega^2 x = 0, \qquad \omega = \sqrt{\frac{k}{m}} .
\end{equation}

\textbf{Allgemeine Lösung} mit Amplitude $A$ und Phasenverschiebung $\varphi$:
\begin{equation}
  x(t) = A\cos(\omega t + \varphi), \qquad
  v(t) = -\omega A \sin(\omega t + \varphi) .
\end{equation}

\textbf{Kreisfrequenz, Frequenz, Periodendauer} klar auseinanderhalten:
\begin{equation}
  \omega = 2\pi f, \qquad T = \frac{2\pi}{\omega} .
\end{equation}

\textbf{Energieerhaltung} im reibungsfreien Fall, nützlich, wenn nach
Geschwindigkeit statt nach Zeit gefragt ist:
\begin{equation}
  E = \tfrac{1}{2} k A^2 = \tfrac{1}{2} m v_{\max}^2 .
\end{equation}
\end{werkzeugkasten}

\begin{werkzeugkasten}[Vorgehen in allgemeiner Form]
Das Muster ist immer dasselbe: Kraftgesetze aufschreiben, einsetzen, auf die
Form $\ddot{x} + \omega^2 x = 0$ bringen und daraus $\omega$ ablesen. Welche
Größe gefragt ist, entscheidet, ob die Zeitfunktion $x(t)$ oder die
Energieerhaltung der kürzere Weg ist. Einheiten zuerst am Ergebnis prüfen, die
Zahlen erst zuletzt einsetzen.
\end{werkzeugkasten}

\begin{achtung}[Typische Fallen]
\begin{itemize}
  \item Gewichtskraft und Federkraft verwechseln. Die Gewichtskraft verschiebt
        nur die Ruhelage, in der Bewegungsgleichung um die Ruhelage fällt sie
        heraus.
  \item Kreisfrequenz $\omega$ und Frequenz $f$ vertauschen, Faktor $2\pi$
        verlieren.
  \item Amplitude $A$ mit der Ruhelage verwechseln und die Verschiebung durch
        die Gewichtskraft zur Auslenkung dazuzählen.
  \item Einheit der Federkonstante falsch verwenden, $\si{\newton\per\metre}$
        ist $\si{\kilo\gram\per\second\squared}$, nicht $\si{\kilo\gram\per\metre}$.
  \item Auslenkung in Zentimetern einsetzen und das Ergebnis in Metern
        erwarten.
\end{itemize}
\end{achtung}
